\documentclass[12pt]{article}

\usepackage[utf8]{inputenc}
\usepackage[T1]{fontenc}
\usepackage[ngerman]{babel}

\usepackage{import}
\usepackage{tabularx}
\usepackage{listings}
\usepackage{hyperref}
\usepackage{CJKutf8}
\usepackage{graphicx}
\usepackage{caption}
\usepackage{microtype}
\usepackage{amssymb}
\usepackage{tikz}
\usetikzlibrary{arrows}
\usetikzlibrary{arrows.meta}
\usepackage{fancyhdr}

\usepackage{times}
\usepackage[margin=2cm,a4paper,headsep=1.2cm]{geometry}

\renewcommand{\baselinestretch}{1.1}
\setlength{\parindent}{0pt}
\setlength{\parskip}{1em}

\setlength{\headheight}{1cm}
\setlength{\headwidth}{\textwidth}

\newcommand{\quotes}[1]{``#1''}

\makeatletter
\newcommand{\labeltext}[2]{%
  \@bsphack
  \csname phantomsection\endcsname % in case hyperref is used
  \def\@currentlabel{#1}{\label{#2}}%
  \@esphack
}
\makeatother    
    
\newcommand{\class}[1]{\textit{#1}-Klasse}
\newcommand{\method}[1]{\textit{#1}-Methode}    
    
\newcounter{fa}
\newcounter{nfa}

\newcommand{\fa}[7]{
\refstepcounter{fa}\label{#7}
\begin{tabularx}{16cm}{l|X}
\textbf{ID} & \textbf{FA\arabic{fa}} \\
\hline
TITEL: & #1 \\
\hline
BESCHREIBUNG: & #2 \\
\hline
BEGRÜNDUNG & #3 \\ 
\hline
ABHÄNGIGKEITEN & #4\\ 
\hline 
PRIORITÄT & #5\\ 
\hline 
AKTEUR & #6
\end{tabularx}
\vspace*{0.75cm}
}

\newcommand{\nfa}[7]{
\refstepcounter{nfa}\label{#7}
\begin{tabularx}{16cm}{l|X}
\textbf{ID} & \textbf{QA\arabic{nfa}} \\
\hline
TITEL: & #1 \\
\hline
BESCHREIBUNG: & #2 \\
\hline
BEGRÜNDUNG & #3 \\ 
\hline
ABHÄNGIGKEITEN & #4\\ 
\hline 
PRIORITÄT & #5\\ 
\hline 
AKTEUR & #6
\end{tabularx}
\vspace*{0.75cm}
}

\newcommand{\akteur}[4]{
\labeltext{#2}{#1}
\begin{tabularx}{16cm}{|l|X|}
%\label{#1}
\hline 
AKTEUR: & #2 \\
\hline
BESCHREIBUNG: & #3 \\
\hline
ROLLE: & #4 \\ 
\hline
\end{tabularx}
}

\newcommand{\myrule}[1]{
	\begin{tikzpicture}
		\draw[{Diamond[open]}-{Diamond[open]}, ultra thick] (0,0) to (#1, 0);
	\end{tikzpicture}

}

\newcommand{\rref}[1]{\ref{#1}\textsuperscript{$\rightarrow$ p. \pageref{#1}}}
\newcommand{\faref}[1]{FA\ref{#1}\textsuperscript{$\rightarrow$ p. \pageref{#1}}}
\newcommand{\mfaref}[1]{FA\ref{#1}}
\newcommand{\nfaref}[1]{QA\ref{#1}\textsuperscript{$\rightarrow$ p. \pageref{#1}}}



\begin{document}


%\maketitle

\begin{titlepage} % Suppresses headers and footers on the title page

	\centering % Centre everything on the title page
	
	\scshape % Use small caps for all text on the title page
	
	\vspace*{\baselineskip} % White space at the top of the page
	
	%------------------------------------------------
	%	Title
	%------------------------------------------------

	
	\myrule{\linewidth}	
	\vspace*{-2\baselineskip}\vspace{4pt}
	\rule{0.95\textwidth}{0.4pt} % Thin horizontal rule
	
%	\includegraphics[width=0.9\textwidth]{images/////Titelblatt}
	
		
	\rule{0.95\textwidth}{0.4pt}
	\vspace*{-2\baselineskip}\vspace{21pt}
	\myrule{\linewidth}
	
	\vspace{2\baselineskip} % Whitespace after the title block
	
	%------------------------------------------------
	%	Subtitle
	%------------------------------------------------
	
	Benutzerhandbuch, \textit{Version: 1.0} - Team 08 \\ 
	Softwaregrundprojekt WiSe 2021/2022 und SoSe 2022
	
	\vspace*{15\baselineskip} % Whitespace under the subtitle

		
	
	%------------------------------------------------
	%	Editor(s)
	%------------------------------------------------
		
	
	bearbeitet von
	
	\vspace{0.5\baselineskip} % Whitespace before the editors
	
	{\scshape Leonhard Alkewitz, Louis Bock, Alexander Fink, \\ Jannis Kieselbach, Felix Schoppe, Niklas Schuerrle\\} % Editor list
	
	\vspace{0.5\baselineskip} % Whitespace below the editor list
	
	\vfill % Whitespace between editor names and publisher logo
	

\end{titlepage}

\tableofcontents
\thispagestyle{empty}

\newpage

\pagestyle{fancy}

\fancyhead[R]{\thepage}
\fancyhead[L]{\leftmark}
\fancyfoot{}


\section{Einführung}

In diesem Projekt geht es um die Entwicklung des Multiplayer-Spiels „Deserts of Dune“.
Das Projekt beinhaltet die Planung als auch Implementierung der vom Stakeholder vorgegebenen
Anforderungen und Wünsche. Die Teammitglieder sollen durch die Umsetzung des Projektes lernen, praktische Erfahrungen mit den Methoden der Lehrveranstaltung zu sammeln. Zum Abschluss des Projektes wird dem Kunden ein funktionstüchtiges Spiel ausgeliefert, welches den Anforderungen entspricht.


\subsection{Motivation}

Die Durchführung des Projektes findet im Rahmen des Informatik-Studiums statt. Die Studenten sollen hierbei das Projekt planen und implementieren. Hierbei sollen die Methoden der Lehrveranstaltung angewandt werden.

\newpage


\section{Der Server}

\subsection{Allgemeines}
Dieses Handbuch behandelt die Handhabung des Servers, welcher von Team 08 im Rahmen des Softwaregrundprojekt im
WISE 2021/2022 und SOSE 2022 an der Universität Ulm entwickelt wurde.

\subsection{Installation und Einrichtung}
Eine Voraussetzung für die Ausführung des Servers ist das Vorliegen der aktuellen Docker Version.

\subsubsection{Bauen und Starten des Docker-Image}
Für den Container liegt ein Dockerfile vor. Mithilfe eines Shell kann die docker-compose.yml Datei ausgeführt werden. Das docker-compose.yml Skript übernimmt das Bauen des Images für den Server. Der Docker-Container kann nun ohne weiteres gestartet werden.

\subsubsection{Kompilieren des Programms}
Es ist auch möglich, die im Ordner .../ befindliche  Server.dll direkt zu starten.

\subsubsection{Kommandozeilenoptionen des Servers}
Der Server besitzt alle im Standard aufgelisteten Konfigurationsmöglichkeiten(auch optionale Kommandozeilenoptionen). Weitere Kommandozeilenoptionen sind nicht vorgesehen.

\textbf{Auflistung der Kommandozeilenoptionen des Servers:}

\textbf{config-party:} Fordert einen absoluten oder relativen Pfad zur .party.json

\textbf{config-scenario:} Fordert einen absoluten oder relativen Pfad zur <Name>.scenario.json

\textbf{x:} Ermöglicht es optionale Schlüssel-Wert-Paare zu übermitteln, die verfügbaren Einstellungen sind von
Server zu Server individuell.

\textbf{help:} Listet alle verfügbaren Parameter auf. Falls individuelle Parameter hinzugefügt wurden so ist dieser Parameter verpflichtend.

\textbf{port:} Fordert nachfolgend eine (gültige) Portnummer, wird keine angegeben oder sollte dieser Parameter
nicht unterstützt werden, so wird die Portnummer 10191 ausgewählt.

\textbf{Auflistung der Kommandozeilenoptionen des KI-Client:}

\textbf{address:} Gibt die IP-Adresse des Servers auf dem gespielt wird an.

\textbf{x:} Ermöglicht es optionale Schlüssel-Wert-Paare zu übermitteln, die verfügbaren Einstellungen sind von
Server zu Server individuell.

\textbf{help:} Listet alle verfügbaren Parameter auf. Falls individuelle Parameter hinzugefügt wurden so ist dieser Parameter verpflichtend.

\textbf{port:} Fordert nachfolgend eine (gültige) Portnummer, wird keine angegeben oder sollte dieser Parameter
nicht unterstützt werden, so wird die Portnummer 10191 ausgewählt.

\textbf{name:} Fordert nachfolgend den Name des Clients im Spiel, wobei dieser mindestens drei Zeichen lang
sein muss.

\end{document}
